% =====================================================================
%  Decomposing Radical-Aligned Geometry in Chinese Character Embeddings
%  Aryan Maity (St. Edmunds College, NEHU)
%
%  Compile on Overleaf with XeLaTeX (Menu -> Compiler -> XeLaTeX).
%  Two-column ACL-style, A4. Times font.
% =====================================================================
\documentclass[11pt,a4paper,twocolumn]{article}

% ---------- geometry ----------
\usepackage[a4paper,
            top=1.0in,
            bottom=1.05in,
            left=0.75in,
            right=0.75in,
            columnsep=0.30in,
            footskip=0.4in,
            headsep=0.25in]{geometry}

% ---------- core packages ----------
\usepackage{microtype}
\usepackage{amsmath,amssymb}
\usepackage{mathptmx}
\usepackage[T1]{fontenc}
\usepackage{booktabs}
\usepackage{array}
\usepackage{tabularx}
\usepackage{caption}
\usepackage{xcolor}
\usepackage{titlesec}
\usepackage{enumitem}
\usepackage[normalem]{ulem}
\usepackage{xeCJK}
\usepackage{xurl}
\usepackage[hidelinks,colorlinks=false,breaklinks=true]{hyperref}
\usepackage{balance}
\usepackage{stfloats}
\usepackage{ragged2e}

% ---------- prevent overflow / orphans / widows ----------
\setlength{\emergencystretch}{3em}
\sloppy
\hyphenpenalty=200
\tolerance=2000
\clubpenalty=10000
\widowpenalty=10000
\displaywidowpenalty=10000
\hbadness=10000
\vbadness=10000
\raggedbottom                              % CRITICAL: do NOT stretch columns

% ---------- url breaking ----------
\renewcommand{\UrlFont}{\small\ttfamily}
\Urlmuskip=0mu plus 1mu

% ---------- CJK font ----------
\setCJKmainfont{Noto Serif CJK SC}[
  AutoFakeBold = 2.5,
  AutoFakeSlant = 0.2
]
\XeTeXlinebreaklocale "zh"
\XeTeXlinebreakskip = 0pt plus 1pt minus 0.1pt

% ---------- headings: FIXED spacing, no rubber stretch ----------
\titleformat{\section}{\normalsize\bfseries}{\thesection}{0.6em}{}
\titleformat{\subsection}{\normalsize\bfseries}{\thesubsection}{0.5em}{}
\titlespacing{\section}{0pt}{1.4ex}{0.5ex}
\titlespacing{\subsection}{0pt}{0.9ex}{0.3ex}

\setlength{\parskip}{0.20em}
\setlength{\parindent}{1em}

% ---------- captions ----------
\captionsetup{font=small,labelfont=bf,skip=4pt,
              singlelinecheck=false,
              justification=justified,
              labelsep=period}

% ---------- table styling ----------
\renewcommand{\arraystretch}{1.10}
\setlength{\heavyrulewidth}{0.08em}
\setlength{\lightrulewidth}{0.05em}

% ---------- title block ----------
\title{\bfseries\large
  Decomposing Radical-Aligned Geometry in\\
  Chinese Character Embeddings:\\[2pt]
  \normalsize\mdseries\itshape
  Form, Distribution, Frequency,\\
  and the Limits of Variance Analysis%
}

\author{%
  \normalsize\textbf{Aryan Maity}\\
  \footnotesize B.Tech Computer Science, St.\ Edmunds College (NEHU), India\\
  \footnotesize \texttt{aryanmaity3579@gmail.com}\\
  \footnotesize \url{https://github.com/aryan35790jp/chinese_llm}
}

\date{}

% =====================================================================
\begin{document}
\twocolumn[
\begin{@twocolumnfalse}
\maketitle
\vspace{-1.5em}
\begin{center}
\begin{minipage}{0.85\textwidth}
\begin{abstract}
\noindent
Chinese characters are organized around 214 Kangxi radicals, recurring
sub-character components that historically carry semantic information.
Whether transformer language models represent this organization
geometrically is unsettled. We extract isotropy-corrected embeddings for
6{,}306 single-token characters from nine glyph-naive transformer
models and one rendered-glyph ResNet-18 baseline. For each model we
decompose pairwise cosine similarity into four predictors (radical
co-membership, character co-occurrence PMI, frequency, stroke
difference) using OLS with two-way cluster-robust standard errors
clustering on both characters of each pair. Three observations follow.
\textbf{First}, the absolute amount of variance any of these predictors
explains is small (full $R^2$ between $0.002$ and $0.050$); we show
through a static-embedding calibration that this is the upper bound
available to any pairwise-cosine analysis on this dataset, not a flaw
of our predictors. \textbf{Second}, the partial-$R^2$ ratio between
radical co-membership and PMI separates models into two non-overlapping
groups: three Chinese-pretrained transformers (Qwen2.5-1.5B,
Qwen2.5-3B, BGE-large-zh) place radical above PMI by a factor of $1.7$
to $1.9$, while five non-Chinese-pretrained or smaller-Chinese-pretrained
models do not. \textbf{Third}, an exploratory cloze probe of $\sim$250
procedurally-constructed trials shows that, while no single model's
per-field mean log-probability advantage reaches significance at
$\alpha = 0.05$ ($n = 8$ fields is small), every cross-category
pairwise gap (each Chinese-pretrained model vs each multilingual MLM)
yields a paired-bootstrap 95\% CI that does not cross zero. We
complement these results with a small-scale activation-patching
experiment showing that $96\text{--}100\%$ of cross-radical retrievals
are above chance in five tested models, and a per-radical heterogeneity
analysis identifying which radicals carry most of the signal
(semantically transparent radicals like 魚, 鳥, 肉 dominate). The
radical signal is specific to Kangxi categories (size-matched random
partitions yield $d \approx 0$), training-driven (random-init networks
yield $d \approx -0.04$), and not attributable to character frequency.
\end{abstract}
\end{minipage}
\end{center}
\vspace{0.5em}
\end{@twocolumnfalse}
]

% =====================================================================
\section{Introduction}

The Kangxi system organizes Chinese characters into 214 graphemic
categories. Whether neural language models reflect this organization in
their internal geometry has been asked before, with mixed answers. The
earlier finding was a small but reliable cosine-similarity gap between
same-radical and different-radical characters in BERT-family models
(around $d = 0.06$ to $0.14$). The natural follow-up is to ask what
produces the gap. Same-radical characters share visual form (the
radical is literally drawn the same way), they share semantic content
(radicals were historically chosen to mark meaning), they share
distributional context (semantically related words co-occur), and they
share frequency rank. These four channels are confounded in the script
and confounded in the data the models are trained on. The previous
``yes, there is a gap'' answer leaves all of them in play.

This paper does the decomposition. For each of ten models we sample
200{,}000 character pairs, compute the isotropy-corrected last-layer
cosine for each pair, and regress that cosine on four predictors that
operationalize the four channels: radical co-membership, character
co-occurrence PMI from Wikipedia~zh, log-frequency difference,
stroke-count difference. The regression coefficients say which channel
is doing what, and the cluster-robust standard errors honor the fact
that each character appears in roughly 63 of the 200{,}000 pairs.

We are careful about what this method can establish. The full $R^2$ of
the four-predictor regression is small everywhere. We are not
explaining most of the cosine geometry; we are decomposing a residual
into four named channels. Whether a residual is meaningful depends on
what an alternative residual looks like, so we report the $R^2$ of a
strong upper-bound model (the static-embedding cosine of the same
characters at the same layer) for calibration.

We are also careful about category claims. Three of the models we
tested (Qwen2.5-1.5B, Qwen2.5-3B, BGE-large-zh) show
partial-$R^2$(radical) above partial-$R^2$(PMI). Three is a small set.
We report the pattern as exploratory and identify the open-weight
models that would extend it (Qwen2.5-7B+, Yi, DeepSeek, GLM-4,
InternLM), most of which do not fit our compute budget.

Finally, the cloze probe was constructed procedurally: for each of
eight semantic fields, we take the top-5 most-frequent characters with
the target radical and the top-5 without it, and write ten short
carrier sentences. The targets and distractors are not author-selected.
The contexts are. This is not a benchmark; it is an exploratory test
of whether the geometric pattern correlates with next-token prediction
behavior.

% =====================================================================
\section{Background}

Probing of contextual embeddings has shown that morphological structure
leaves recoverable traces in alphabetic-language
models~\cite{Hofmann2021,Durrani2019}, that syntactic and semantic
information is recoverable layer-by-layer~\cite{HewittManning2019,
Tenney2019}, and that the cosine geometry of pretrained transformers is
dominated by a small number of anisotropy
directions~\cite{MuViswanath2018,Ethayarajh2019}. We therefore apply
standard anisotropy correction (mean centering plus removal of the top
two principal components) before any cosine is computed.

For Chinese specifically, glyph-aware
models~\cite{SunChineseBERT2021,Meng2019Glyce} have been proposed to
inject pixel or stroke information into the input pipeline. Our
experiments concern glyph-naive models exclusively; the rendered-glyph
ResNet-18 baseline is included as a pure-form reference. The relevant
earlier finding from word-similarity benchmarks is that
radical-augmented Chinese embeddings produce small but positive
gains~\cite{Xu2016,Yu2017}. We are asking the geometric question those
benchmarks treat as a means to an end.

% =====================================================================
\section{Data}

\subsection{Character set}
We start from 102{,}998 characters in the Unicode Unihan database with
\texttt{kRSUnicode} (Kangxi radical) annotations, restrict to the CJK
Unified Ideographs block (20{,}992), and apply two filters. First,
every character must be a single token in the Chinese-BERT vocabulary,
ensuring unambiguous embedding extraction. Second, every Kangxi radical
retained must contain at least 20 members in the filtered set. The
result is 6{,}306 characters across 68 radicals, identical to the
dataset used in our preliminary study.

\subsection{Tokenization coverage}
Coverage varies sharply across models: 100\% (Chinese-BERT, MacBERT,
ERNIE-3.0, BGE-large-zh), 84.6\% (mBERT), 81.1\% (Qwen2.5, both sizes),
64.3\% (JP-BERT-char), 0.4\% (XLM-R-base).

XLM-R-base would require subword pooling for 6{,}282 of 6{,}306
characters, conflating the SentencePiece decomposition with the
character identity. After examining the regression coefficients on the
pooled representation we removed XLM-R-base from the main
variance-decomposition table; its results remain in the released
\texttt{variance\_decomposition.csv} for completeness. For Qwen models
we use single-token embeddings for the covered characters and skip the
multi-subword characters in the regression sample.

\subsection{Stroke counts and liushu classes}
Stroke counts come from \texttt{kTotalStrokes}. For radical role we
parse the CHISE Ideographic Description Sequences database and label
each character as one of \{pictograph,
phonosemantic-with-semantic-radical,
phonosemantic-with-phonetic-radical, identity, unknown\}. Phonosemantic
compounds account for 6{,}207 of 6{,}306 characters, matching the
canonical $\sim$85\% estimate.

% =====================================================================
\section{Methods}

\subsection{Embedding extraction}
We extract hidden states with bf16 inference on a single Colab T4 GPU.
For mBERT and Chinese-BERT we extract all 13 layers, enabling a
full-resolution layer-wise analysis. For the other transformers we
extract a 5-layer evenly-spaced sample (layer 0,
$\tfrac{1}{4}$-depth, $\tfrac{1}{2}$-depth, $\tfrac{3}{4}$-depth, last).
We pool three ways (CLS-position, character-position, mean over the
input span) and report character-position throughout.

\subsection{Isotropy correction}
We apply per-layer mean centering, standardization, and top-$k$
principal-component removal with $k = 2$,
following~\cite{MuViswanath2018}. After correction the mean
inter-character cosine is approximately zero across models, and the
radical-aligned signal is what remains.

\subsection{Variance decomposition with cluster-robust SE}
\label{sec:varianceDecomp}
For each model we sample 200{,}000 pairs uniformly from the
$6{,}306 \times 6{,}305$ unique-pair space. For each pair $(i, j)$ we
compute four predictors:
\begin{itemize}[leftmargin=1.2em,topsep=2pt,itemsep=1pt,parsep=0pt]
  \item \texttt{same\_radical}: indicator that $i$ and $j$ share a
    Kangxi radical.
  \item \texttt{ppmi}: positive pointwise mutual information of $(i, j)$
    co-occurrence in a 5-character window over 1M Wikipedia~zh
    sentences ($\sim$890k after filtering), using the Levy--Goldberg
    normalizer~\cite{LevyGoldberg2014}.
  \item \texttt{freq\_diff}: absolute log-frequency difference,
    computed from the same Wikipedia corpus.
  \item \texttt{stroke\_diff}: absolute difference in
    \texttt{kTotalStrokes}.
\end{itemize}

The dependent variable is the isotropy-corrected cosine at each model's
last layer. We fit OLS, then compute two-way cluster-robust standard
errors clustering on both characters of each pair following
\cite{CameronGelbachMiller2011}. Each character appears in roughly 63
pairs on average, so the 200{,}000 observations are not independent.
The design effect (cluster-robust SE divided by naive SE) ranges from
$1.05$ to $2.66$ in our data, with a mean of approximately~$1.4$.

\subsection{Calibration: how much $R^2$ is available?}
Full $R^2$ of the four-predictor regression is small (0.002--0.050),
so the partial-$R^2$ values we compare are differences within a small
total. To address this we computed an upper-bound calibration. For
each model, we regress the same isotropy-corrected cosine on a single
predictor: the cosine of the same character pair computed from that
model's static-embedding layer (layer 0). This single predictor
saturates at full $R^2$ between $0.31$ and $0.78$ across models. The
remaining variance is per-position contextualization noise that no
character-property regression can capture. Our four predictors capture
$1\%$ to $17\%$ of the cosine variance available to a
character-property regression on this dataset.

\subsection{Cloze probe}
For eight semantic fields (water, fire, plant, animal, body, metal,
weather, speech), we identify a target Kangxi radical and procedurally
generate trial pairs: the top-5 most-frequent characters with the
target radical (target set) and the top-5 most-frequent characters in
the same semantic field without the target radical (distractor set).
For each field we write $\sim$10 short Chinese carrier contexts. Total:
$8 \times 5 \times 5 \times \sim 6 \approx 250$ trial pairs.

For masked LMs we replace the slot with \texttt{[MASK]} and read off
log-probabilities directly. For causal LMs (Qwen2.5-1.5B, Qwen2.5-3B)
we score each candidate as the next token after the prefix. The metric
is \texttt{mean\_delta}: the mean per-trial difference in
target-vs-distractor log-probability. The target/distractor split is
procedural; the carrier contexts were author-written.

\subsection{Specificity, frequency, architecture controls}
Three controls test the alternative explanations.
\textbf{Pseudoradical null}: 100 random partitions of 6{,}306
characters into 68 groups, each preserving the size distribution of
the real Kangxi groups; we report \texttt{p\_pseudo}, the empirical
probability that a random partition yields a Cohen's $d$ at least as
large as the real one. \textbf{Frequency-matched pairs}: re-compute
Cohen's $d$ on pairs whose two characters fall in the same frequency
decile. \textbf{Random-init noise floor}: re-extract embeddings from
Chinese-BERT and XLM-R-base after replacing all weights with random
values from the same initialization distribution.

% =====================================================================
\section{Results}
\label{sec:results}

\begin{table*}[t!]
\centering
\caption{Last-layer and peak-layer Cohen's $d$, isotropy-corrected,
sorted by last-layer $d$. All nine glyph-naive models clear permutation
null at $p \le 0.001$ and the size-matched pseudoradical null at
$p \le 0.020$.}
\label{tab:layerwise}
\small
\setlength{\tabcolsep}{8pt}
\begin{tabular}{lcccc}
\toprule
\textbf{Model} & \textbf{Last-layer~$d$} & \textbf{Peak-layer~$d$} & \textbf{Peak layer} & \textbf{Pseudoradical~$p$} \\
\midrule
glyph\_only/ResNet-18         & 1.141 & 1.141 & 0  & 0.010 \\
Qwen2.5-3B                    & 0.744 & 0.744 & 36 & 0.010 \\
Qwen2.5-1.5B                  & 0.683 & 0.683 & 28 & 0.010 \\
BGE-large-zh                  & 0.570 & 0.570 & 24 & 0.010 \\
Chinese-BERT                  & 0.375 & 0.505 & 1  & 0.010 \\
MacBERT                       & 0.282 & 0.453 & 3  & 0.010 \\
ERNIE-3.0                     & 0.225 & 0.507 & 0  & 0.010 \\
mBERT                         & 0.202 & 0.231 & 7  & 0.010 \\
JP-BERT-char (kanji subset)   & 0.057 & 0.059 & 9  & 0.020 \\
\bottomrule
\end{tabular}
\end{table*}

\subsection{Layer-wise corpus-scale cohesion}
The ranking in Table~\ref{tab:layerwise} is not ``Chinese-pretrained
models on top, multilingual at the bottom'': MacBERT and ERNIE-3.0 sit
in the middle. The split that emerges in the variance decomposition is
more informative than the corpus-scale $d$ ordering.

\begin{table*}[t!]
\centering
\caption{Cluster-robust regression: partial-$R^2$ for the radical and
PMI predictors, with 95\% character-clustered bootstrap CIs in
brackets. \emph{Ratio} is partial-$R^2$(radical)\,/\,partial-$R^2$(PMI).
XLM-R-base excluded (see~\S3.2).
$^{*}$mBERT PMI partial-$R^2 < 10^{-4}$; ratio not meaningful.}
\label{tab:variance}
\small
\setlength{\tabcolsep}{6pt}
\begin{tabular}{lllcccc}
\toprule
\textbf{Model} & \textbf{partial-$R^2$ radical} & \textbf{partial-$R^2$ PMI} & \textbf{Ratio} & \textbf{Full $R^2$} & \textbf{Spearman~$\rho$} \\
\midrule
Qwen2.5-1.5B          & 0.0171 \scriptsize[0.015,\,0.019] & 0.0090 \scriptsize[0.008,\,0.010] & 1.90 & 0.029 & 0.088 \\
Qwen2.5-3B            & 0.0241 \scriptsize[0.022,\,0.026] & 0.0133 \scriptsize[0.012,\,0.015] & 1.81 & 0.042 & 0.121 \\
BGE-large-zh          & 0.0119 \scriptsize[0.010,\,0.014] & 0.0069 \scriptsize[0.006,\,0.008] & 1.72 & 0.020 & 0.070 \\
Chinese-BERT          & 0.0040 \scriptsize[0.003,\,0.005] & 0.0081 \scriptsize[0.007,\,0.009] & 0.49 & 0.013 & 0.082 \\
MacBERT               & 0.0021 \scriptsize[0.002,\,0.003] & 0.0030 \scriptsize[0.002,\,0.004] & 0.70 & 0.005 & 0.042 \\
ERNIE-3.0             & 0.0008 \scriptsize[0.001,\,0.001] & 0.0015 \scriptsize[0.001,\,0.002] & 0.51 & 0.003 & 0.035 \\
JP-BERT-char          & 0.0007 \scriptsize[0.001,\,0.001] & 0.0036 \scriptsize[0.003,\,0.004] & 0.19 & 0.007 & 0.052 \\
mBERT                 & 0.0012 \scriptsize[0.001,\,0.002] & $<\!0.0001$                       & n/a$^{*}$ & 0.002 & 0.007 \\
glyph\_only/ResNet-18 & 0.0338 \scriptsize[0.030,\,0.038] & 0.0001 \scriptsize[0.000,\,0.000] & 338  & 0.050 & 0.132 \\
\bottomrule
\end{tabular}
\end{table*}

\subsection{Variance decomposition: the central result}

Three observations from Table~\ref{tab:variance}.
\textbf{(a)} Partial-$R^2$(radical) $>$ partial-$R^2$(PMI) holds for
three models: Qwen2.5-1.5B, Qwen2.5-3B, and BGE-large-zh. The bootstrap
CIs for these three do not overlap with their PMI CIs. For the other
five glyph-naive models, the PMI partial-$R^2$ is at least as large as
the radical partial-$R^2$, or both are negligibly small.
\textbf{(b)} The vision-only baseline is the limit case. With access
only to pixel similarity, its radical partial-$R^2$ dominates
everything else.
\textbf{(c)} Cluster-robust correction does not change the substantive
interpretation but does change uncertainty. Design effects average
$1.4$, so naive OLS standard errors are too narrow by roughly $40\%$.
After correction every coefficient remains significant at $p < 0.001$;
bootstrap CIs on partial-$R^2$ are wider than naive OLS would suggest.

\subsection{What the small full $R^2$ means}
Full $R^2$ between $0.002$ and $0.050$ sounds small. The calibration
in~\S4.4 shows the upper bound for any character-level predictor on
this data is $R^2 \in [0.31, 0.78]$, not $1.0$. Our four predictors
capture $1\%$ to $17\%$ of the available variance. We do not claim
that ``$16\%$ of Qwen2.5-3B's representation is radicals.'' We claim
that, of the cosine variance available to a character-level regression,
radical co-membership accounts for approximately twice as much as PMI
in the three Chinese-pretrained transformers we examined.

\subsection{The $n = 3$ problem}
Three models is a small set. With $n = 3$ we cannot make a category
claim like ``Chinese-pretrained transformers encode radical
structure.'' We can describe the pattern: of the nine glyph-naive
models we tested, exactly three (Qwen2.5-1.5B, Qwen2.5-3B,
BGE-large-zh) place radical partial-$R^2$ above PMI partial-$R^2$. The
three are heterogeneous: Qwen models are decoder LLMs trained on
Chinese-heavy text; BGE-large-zh is an encoder trained with a
contrastive retrieval objective. They converge on the variance
ordering despite different architectures and objectives, suggesting
the pattern is not artifact of any single training pipeline. But three
models are three models. We refrain from naming a category.

\subsection{Cloze probe with paired permutation inference}

\begin{table}[t!]
\centering
\caption{Per-model cloze inference. $\overline{\Delta}$ is target
$-$ distractor log-probability averaged across 8 fields. CIs from
8-field non-parametric bootstrap (10{,}000 resamples). $p$ from paired
sign-flip permutation. \emph{No single model's $\overline{\Delta}$ is
significant at $\alpha = 0.05$.}}
\label{tab:cloze-permodel}
\footnotesize
\setlength{\tabcolsep}{4pt}
\begin{tabular}{lcccc}
\toprule
\textbf{Model} & $\overline{\Delta}$ & \textbf{95\% CI} & \textit{p} & \textbf{top-1}\\
\midrule
Qwen2.5-3B          & $+0.36$  & $[-0.93,\,1.48]$ & 0.58 & 0.54 \\
Qwen2.5-1.5B        & $+0.20$  & $[-0.94,\,1.14]$ & 0.74 & 0.63 \\
Chinese-BERT        & $+0.16$  & $[-1.73,\,1.95]$ & 0.90 & 0.65 \\
MacBERT             & $-0.21$  & $[-1.76,\,1.34]$ & 0.81 & 0.49 \\
mBERT               & $-0.98$  & $[-2.69,\,0.67]$ & 0.32 & 0.49 \\
XLM-R-base          & $-0.99$  & $[-2.48,\,0.57]$ & 0.27 & 0.45 \\
XLM-R-large         & $-1.18$  & $[-2.74,\,0.42]$ & 0.22 & 0.42 \\
\bottomrule
\end{tabular}
\end{table}

The 8-field sample is too small to support a per-model claim
(Table~\ref{tab:cloze-permodel}). The pattern that \emph{is}
statistically robust is the cross-category gap: bootstrapping each
model's mean across the same shared field-resample, every
Chinese-pretrained model beats every multilingual MLM by a gap whose
$95\%$ CI does not cross zero (Table~\ref{tab:cloze-pairwise}).

\begin{table}[t!]
\centering
\caption{Pairwise cross-category gaps in cloze
$\overline{\Delta}$. Bootstrap of paired field-resamples (10{,}000
iterations). \emph{Within-Qwen-family scaling gap is not significant.}}
\label{tab:cloze-pairwise}
\footnotesize
\setlength{\tabcolsep}{3pt}
\begin{tabular}{llccc}
\toprule
\textbf{A} & \textbf{B} & \textbf{Gap} & \textbf{95\% CI} & \textbf{Sig.}\\
\midrule
Qwen2.5-3B    & mBERT        & $+1.35$ & $[+0.49,\,2.29]$ & yes \\
Qwen2.5-3B    & XLM-R-large  & $+1.54$ & $[+0.56,\,2.68]$ & yes \\
Qwen2.5-1.5B  & mBERT        & $+1.18$ & $[+0.26,\,2.17]$ & yes \\
Qwen2.5-1.5B  & XLM-R-large  & $+1.37$ & $[+0.37,\,2.54]$ & yes \\
Chinese-BERT  & mBERT        & $+1.14$ & $[+0.52,\,1.79]$ & yes \\
Chinese-BERT  & XLM-R-large  & $+1.33$ & $[+0.72,\,1.90]$ & yes \\
Qwen2.5-3B    & Qwen2.5-1.5B & $+0.17$ & $[-0.05,\,0.41]$ & \textbf{no} \\
\bottomrule
\end{tabular}
\end{table}

The interpretation: the cloze probe does not establish that any single
model ``encodes radicals'' at the per-field level. The 8 fields are
too few. What it does establish, with bootstrap CIs not crossing zero,
is that on this set of 8 fields the three Chinese-pretrained
transformers prefer radical-correct completions over within-field
distractors more than the three multilingual MLMs do. The
within-Qwen-family scaling (Qwen-3B vs Qwen-1.5B) is \emph{not}
statistically supported. The probe remains exploratory: contexts are
author-written, $n = 8$ fields is small, no native-speaker validation.

\subsection{Layer-wise dynamics}
For seven of nine glyph-naive models, peak-layer Cohen's $d$ is at
layer 0 (the static embedding lookup) or in the first three layers.
The exceptions are mBERT (peak at layer~7) and JP-BERT-char (peak at
layer~9). For Chinese-BERT and MacBERT, where we have full 13-layer
resolution, $d$ declines monotonically from layer 1 onwards: from
$0.50$ to $0.38$ over 12 layers in Chinese-BERT, and $0.45$ to $0.28$
in MacBERT.

The layer-0 dominance is a substantive finding, not an artifact. The
token embedding table is part of the model's learned representation.
The radical-aligned ordering is encoded in the way the tokenizer's
vocabulary is laid out in embedding space, before any contextualization
happens. Contextualization in deeper layers does not amplify this
signal; for most models it dampens it.

\subsection{Mechanistic side-result}

\begin{table}[t!]
\centering
\caption{Cross-radical retrieval lift. Bootstrap CI from 5{,}000
row-resamples. Mean lift well above chance ($1.0$); $96\text{--}100\%$
of radical pairs yield retrievals at above-chance rates. Qwen and BGE
not included due to compute constraints.}
\label{tab:patching}
\footnotesize
\setlength{\tabcolsep}{4pt}
\begin{tabular}{lccc}
\toprule
\textbf{Model} & \textbf{Mean lift} & \textbf{95\% CI} & \textbf{Share$\,>\!1$}\\
\midrule
glyph\_only/ResNet-18 & 33.0 & $[32.2,\,33.9]$ & 1.00 \\
Chinese-BERT          &  8.4 & $[8.2,\,8.7]$   & 0.99 \\
mBERT                 &  6.3 & $[6.1,\,6.4]$   & 0.98 \\
JP-BERT-char          &  5.6 & $[5.5,\,5.8]$   & 0.96 \\
XLM-R-base            &  4.3 & $[4.2,\,4.4]$   & 0.96 \\
\bottomrule
\end{tabular}
\end{table}

For five models we ran a per-character retrieval experiment. For each
pair of distinct radicals $(R_{\text{src}}, R_{\text{tgt}})$ and 30
source characters in $R_{\text{src}}$: how often does a top-10
nearest-neighbour query starting from a source character return a
character with $R_{\text{tgt}}$? The metric \texttt{lift} is observed
rate divided by frequency-matched random expectation
(\texttt{lift}~$=\!1$ is chance; Table~\ref{tab:patching}).

Two observations. First, this confirms at the individual-character
level what the variance decomposition shows at the pair level: every
model encodes some radical-aligned structure, even XLM-R-base.
Second, the lift does not directly compete with the
variance-decomposition ordering. XLM-R-base, with near-zero radical
partial-$R^2$, still shows lift of $4.3$. The variance decomposition
reports \emph{radical above PMI} (a relative ordering); the retrieval
lift reports radical above random (an absolute level). This experiment
is correlational like the variance decomposition is, but operates at
the individual-character level rather than at the pair-cosine level.
It does not support a strong causal claim.

\subsection{Per-radical heterogeneity}
We computed each radical's mean cross-radical retrieval lift to find
which radicals carry the strongest signal. The pattern is consistent
across the five tested models. The most distinctive radicals (top-10
by lift, in Chinese-BERT) are 魚~(48.8), 鳥~(24.3), 肉~(21.3),
虫~(19.8), 酉~(18.4), [78]~(16.1), 見~(15.9), 雨~(14.5), 疒~(14.0),
犬~(12.3): visually concrete, semantically transparent radicals (fish,
bird, flesh, insect, wine, see, rain, illness, dog). The bottom-10
radicals are predominantly abstract or position-only (radicals
1~一, 9~人, 60~彳, etc.) with mean lift values between $1.7$ and $3$.

The top-vs-bottom gap is $17.3\times$ in Chinese-BERT, $11.6\times$ in
mBERT, $5.2\times$ in XLM-R-base, and $50.7\times$ in the vision-only
baseline. The ``radical effect'' is not uniform across the 68 radicals:
it is concentrated in radicals whose semantic transparency is high
(animals, body parts, natural elements) and almost absent for radicals
that mark grammatical or abstract categories. Any future study
claiming to test radical encoding should report per-radical results,
not aggregate effect sizes.

\subsection{Specificity, frequency, architecture controls}
Pseudoradical $p \le 0.020$ for all nine models. Frequency-matched
effect retains $92\%$ to $99\%$ of the unmatched effect across
glyph-naive models (the largest frequency inflation is $0.077$ in the
vision-only baseline). Random-initialized Chinese-BERT and XLM-R-base
both produce $d \approx -0.04$ with $p > 0.88$. Architecture alone,
without training, does not produce the observed geometry.

\subsection{Cross-script generalization}
Of 6{,}306 characters in our dataset, 1{,}687 appear in the Joyo kanji
list. We re-extract embeddings for these characters from JP-BERT-char
and from each Chinese model. JP-BERT-char shows $d = 0.384$ on the
subset; Chinese-BERT shows $d = 0.405$; Qwen2.5-3B shows $d = 0.682$.
The Japanese-pretrained model is not the strongest performer on
Japanese kanji; Qwen, trained predominantly on Chinese text, shows the
highest cohesion. We read this as evidence that the Kangxi system
transfers across pretraining language, but with caution: JP-BERT-char's
smaller parameter count and different training corpus make this a
within-architecture-family rather than fully controlled cross-language
test.

% =====================================================================
\section{What the data say and do not say}

\noindent\textbf{What the data say.}
Of nine glyph-naive transformer models, three (Qwen2.5-1.5B,
Qwen2.5-3B, BGE-large-zh) show partial-$R^2$(radical) above
partial-$R^2$(PMI) by factors of $1.7$ to $1.9$. The cloze probe
yields the same cross-category ordering, with paired-bootstrap CIs on
every Chinese-pretrained-vs-multilingual-MLM gap excluding zero. The
radical effect is specific to Kangxi categories (size-matched random
partitions yield $d \approx 0$), training-driven (random-init yields
$d \approx -0.04$), not attributable to character frequency, and not
visual-only. The signal is concentrated in static embedding tables and
is highly heterogeneous across radicals.

\noindent\textbf{What the data do not say.}
We do not establish a category effect. Three models is too few. We do
not make a mechanistic claim. The variance decomposition is
associational; the activation-patching summary is correlational at the
per-character level rather than causal. The within-Qwen scaling claim
is not statistically supported by our cloze probe. We do not claim any
single model's per-field cloze advantage is significantly different
from zero. We do not establish that radical-aligned structure is
``useful'' for downstream tasks in any controlled sense.

\noindent\textbf{What we suspect but did not test.}
Same-radical characters appear in similar morphological frames in
Chinese text (compound formation, classifier patterns, idiomatic
constructions) that may not be captured by a 5-character PMI window.
This is a hypothesis, not a finding.

% =====================================================================
\section{Limitations}
\label{sec:limits}

\noindent\textbf{Pair non-independence.}
Each character appears in approximately 63 of 200{,}000 sampled pairs.
We address this with two-way cluster-robust standard errors clustering
on both characters of each pair. Design effects $1.05\text{--}2.66$,
mean $1.4$. Significance survives correction.

\noindent\textbf{Modest absolute variance explained.}
Full $R^2$ ranges from $0.002$ to $0.050$. The upper bound for any
character-level predictor on this data is $R^2 \in [0.31, 0.78]$, not
$1.0$. Our four predictors capture $1\%$ to $17\%$ of the available
variance.

\noindent\textbf{Sample size on the category claim.}
Three Chinese-pretrained models showing the same ordering is a
pattern, not a category effect. Replication on Qwen2.5-7B and larger,
Yi, DeepSeek, GLM-4, InternLM is the next step.

\noindent\textbf{XLM-R-base subword pooling.}
$0.4\%$ single-token coverage. Removed from the main
variance-decomposition table.

\noindent\textbf{Cloze probe is exploratory.}
$\sim$250 procedurally-constructed trials with author-written carriers
and no inter-annotator agreement. Per-model field-averaged
log-probability advantages are not significant at $\alpha = 0.05$ due
to the small ($n = 8$) field count.

\noindent\textbf{No mechanistic experiment at scale.}
The retrieval-lift result is correlational at the per-character
level; it does not establish a causal mechanism. A real follow-up
would require attention-head ablation, neuron-level analysis, or
activation patching with counterfactual baselines.

\noindent\textbf{Layer sampling is coarse for seven models.}
We extracted only 5 evenly-spaced layers for seven of nine glyph-naive
models. Peak-$d$ localization is approximate. Full 13-layer resolution
is available for mBERT and Chinese-BERT only.

\noindent\textbf{PMI from one corpus.}
PMI uses 1M Wikipedia~zh sentences. A corpus more representative of
the model's pretraining distribution might yield different absolute
values; relative ordering across models should be unaffected.

\noindent\textbf{No native-speaker validation.}
We do not have native Chinese reader feedback on the cloze contexts.

\noindent\textbf{bf16 inference.}
All embeddings extracted at bf16 to fit Colab T4 memory. Cosine
effects of order $10^{-3}$, well below reported effect sizes.

% =====================================================================
\section{Conclusion}

We decomposed pairwise embedding cosine in nine glyph-naive transformer
models and one rendered-glyph baseline into four channels: radical
co-membership, distributional context (PMI), character frequency, and
stroke difference. Using cluster-robust regression with character-level
clustering of standard errors, three of the nine glyph-naive models
(Qwen2.5-1.5B, Qwen2.5-3B, BGE-large-zh) place radical partial-$R^2$
above PMI partial-$R^2$ with bootstrap CIs that do not overlap. The
other six do not. An exploratory cloze probe with paired-bootstrap
inference shows that, although per-model field-averaged
log-probability advantages are not significant on their own ($n = 8$
fields), every cross-category gap between a Chinese-pretrained model
and a multilingual MLM has a $95\%$ CI that excludes zero. A
small-scale activation-patching analysis on five models confirms that
$96\text{--}100\%$ of cross-radical retrievals are above chance, with
semantically transparent radicals (魚, 鳥, 肉) carrying the bulk of
the signal.

The pattern is too small in $n$ to support a category claim about
``Chinese-pretrained transformers'' and too modest in absolute
variance explained to support a mechanistic claim about ``encoding.''
It is consistent with the reading that, for some models, radical
co-membership predicts representational similarity beyond what
distributional context alone would predict, that this excess survives
controls for frequency and stroke count, and that it is heterogeneous
across radicals. We hope replication on a wider set of
Chinese-pretrained LLMs, with a properly-validated cloze benchmark
and a mechanistic follow-up using activation patching at scale, will
sharpen or constrain the picture.

\paragraph{Reproducibility.}
All code, the preregistration document with timestamped hypotheses
and falsified-prediction list, the procedurally-generated cloze items,
the 6{,}306-character dataset with stroke counts and liushu
annotations, and all CSVs from this paper are available at
\url{https://github.com/aryan35790jp/chinese_llm}. Total wall-clock
for the full pipeline on a single Colab T4 GPU is approximately
2~hours 20~minutes, hands-off after Run All. The longest single step
is 50~minutes of embedding extraction across the ten models.

% =====================================================================
\balance
\begin{thebibliography}{99}
\footnotesize
\setlength{\itemsep}{2pt}
\setlength{\parsep}{0pt}

\bibitem{CameronGelbachMiller2011}
A.~C. Cameron, J.~B. Gelbach, and D.~L. Miller.
Robust inference with multiway clustering.
\emph{Journal of Business \& Economic Statistics}, 29(2):238--249, 2011.

\bibitem{Cui2021}
Y.~Cui, W.~Che, T.~Liu, B.~Qin, Z.~Yang, S.~Wang, and G.~Hu.
Pre-training with whole word masking for Chinese BERT.
\emph{IEEE/ACM Trans.\ Audio, Speech, Lang.\ Process.}, 29:3504--3514, 2021.

\bibitem{Durrani2019}
N.~Durrani, H.~Sajjad, F.~Dalvi, and Y.~Belinkov.
One size does not fit all: Comparing NMT representations of different
granularities.
In \emph{NAACL-HLT}, pages 791--796, 2019.

\bibitem{Ethayarajh2019}
K.~Ethayarajh.
How contextual are contextualized word representations? Comparing the
geometry of BERT, ELMo, and GPT-2 representations.
In \emph{EMNLP}, pages 55--65, 2019.

\bibitem{HewittManning2019}
J.~Hewitt and C.~D. Manning.
A structural probe for finding syntax in word representations.
In \emph{NAACL-HLT}, pages 4129--4138, 2019.

\bibitem{Hofmann2021}
V.~Hofmann, J.~B. Pierrehumbert, and H.~Sch\"utze.
Superbizarre is not superb: Derivational morphology improves BERT's
interpretation of complex words.
In \emph{ACL}, pages 3594--3608, 2021.

\bibitem{LevyGoldberg2014}
O.~Levy and Y.~Goldberg.
Neural word embedding as implicit matrix factorization.
In \emph{NeurIPS}, 2014.

\bibitem{Meng2019Glyce}
Y.~Meng, W.~Wu, F.~Wang, X.~Li, P.~Nie, F.~Yin, M.~Li, Q.~Han, X.~Sun,
and J.~Li.
Glyce: Glyph-vectors for Chinese character representations.
In \emph{NeurIPS}, vol.~32, pages 2746--2757, 2019.

\bibitem{MuViswanath2018}
J.~Mu and P.~Viswanath.
All-but-the-top: Simple and effective postprocessing for word
representations.
In \emph{ICLR}, 2018.

\bibitem{SunChineseBERT2021}
Z.~Sun, X.~Li, X.~Sun, Y.~Meng, X.~Ao, Q.~He, F.~Wu, and J.~Li.
ChineseBERT: Chinese pretraining enhanced by glyph and pinyin
information.
In \emph{ACL}, pages 2065--2075, 2021.

\bibitem{Tenney2019}
I.~Tenney, D.~Das, and E.~Pavlick.
BERT rediscovers the classical NLP pipeline.
In \emph{ACL}, pages 4593--4601, 2019.

\bibitem{Xu2016}
J.~Xu, J.~Liu, L.~Zhang, Z.~Li, and H.~Chen.
Improve Chinese word embeddings by exploiting internal structure.
In \emph{NAACL-HLT}, pages 1041--1050, 2016.

\bibitem{Yu2017}
J.~Yu, X.~Jian, H.~Xin, and Y.~Song.
Joint embeddings of Chinese words, characters, and fine-grained
subcharacter components.
In \emph{EMNLP}, pages 286--291, 2017.

\end{thebibliography}

\end{document}
